\documentclass[11pt]{article}

\usepackage[letterpaper,margin=1in]{geometry}
\usepackage{graphicx}
\usepackage{booktabs}
\usepackage{tabularx}
\usepackage{array}
\usepackage{float}
\usepackage{xcolor}
\usepackage{hyperref}
\usepackage{enumitem}

\hypersetup{
    colorlinks=true,
    linkcolor=blue!50!black,
    urlcolor=blue!60!black
}

\newcommand{\projecttitle}{OFT Finetuning for Schema-Grounded Text-to-SQL}
\newcommand{\coursename}{AIST5030}
\newcommand{\teamname}{Mini-Project Team}
\newcommand{\memberlist}{Member 1, Member 2, Member 3}
\newcommand{\modelname}{Qwen/Qwen2.5-0.5B-Instruct}
\newcommand{\datasetname}{b-mc2/sql-create-context}

\newcommand{\insertfigure}[3]{
    \begin{figure}[H]
        \centering
        \IfFileExists{#1}{
            \includegraphics[width=#2\linewidth]{#1}
        }{
            \fbox{
                \begin{minipage}[c][0.25\textheight][c]{0.92\linewidth}
                    \centering
                    Missing figure file: \texttt{\detokenize{#1}}\\[0.4em]
                    Replace with your generated plot before submission.
                \end{minipage}
            }
        }
        \caption{#3}
    \end{figure}
}

\begin{document}

\begin{center}
    {\LARGE \textbf{\projecttitle}}\\[0.5em]
    {\large \coursename\ Mini-Project Report Template}\\[0.5em]
    \teamname\\
    \memberlist\\
    \today
\end{center}

\vspace{0.5em}
\noindent\textbf{How to use this template.}
This version is pre-filled using the actual outputs in \texttt{outputs/train\_4h} and \texttt{outputs/eval}.
This report can be compiled directly; edit team metadata if needed and optionally expand qualitative examples.

\section{Task and Setup}
\textbf{Task.}
Finetune a pretrained LLM with OFT for schema-grounded Text-to-SQL generation.

\textbf{Input/Output.}
Input is a natural language question with SQL schema context (CREATE TABLE statements);
output is a SQL query.

\begin{table}[H]
    \centering
    \caption{Reproducibility Configuration}
    \begin{tabularx}{\linewidth}{>{\raggedright\arraybackslash}p{0.28\linewidth}X}
        \toprule
        Item & Value \\
        \midrule
        Base model & \modelname \\
        Dataset & \datasetname \\
        Train/Val/Test split sizes & 20000 / 2000 / 2000 \\
        OFT target modules & \texttt{q\_proj, k\_proj, v\_proj, o\_proj, up\_proj, down\_proj, gate\_proj} \\
        Max sequence length & 384 \\
        Training budget & Single-GPU bounded run (\(\sim\)3.17 hours train runtime; 1800 max steps) \\
        Key files & \texttt{configs/sql\_4h.yaml}, \texttt{src/train\_oft\_sql.py}, \texttt{src/eval\_sql.py} \\
        \bottomrule
    \end{tabularx}
\end{table}

\section{Method}
\textbf{Prompting.}
Each example includes a natural-language question plus schema context (CREATE TABLE statements),
and the model is prompted to generate SQL only. During evaluation, outputs are normalized to lower-case,
compacted whitespace, and stripped trailing semicolons before exact-match comparison.

\textbf{OFT Design.}
We attach OFT adapters to both attention projections (Q/K/V/O) and MLP projections
(up/down/gate) to adapt both schema-linking behavior and compositional reasoning while keeping
the base model architecture fixed.

\textbf{Training Strategy.}
Training uses per-device batch size 2 with gradient accumulation 16, cosine LR schedule,
learning rate \(2\times10^{-4}\), warmup ratio 0.05, 2 epochs with max 1800 steps,
evaluation/checkpoint every 100 steps, and early-stopping patience of 3.

\section{Results}
\subsection{Training Dynamics}
\insertfigure{../outputs/plots/loss_curves.png}{0.9}{Training and validation loss curves.}
\noindent Training was stable across the bounded run, with final train loss of 0.644 and test loss of 0.616.
The low gap between train and test loss suggests no strong overfitting under this budget.
The run completed in 11427 seconds (\(\sim\)3.17 hours), matching the intended 4-hour envelope.

\subsection{Baseline vs OFT Metrics}
\begin{table}[H]
    \centering
    \caption{Evaluation Metrics on the Fixed Test Split}
    \begin{tabularx}{\linewidth}{lcc>{\raggedright\arraybackslash}X}
        \toprule
        Model & Exact Match (\%) & SQL Parse Success (\%) & Notes \\
        \midrule
        Baseline (pre-finetune) & 0.15 & 100.00 & Mostly syntactically valid SQL but weak schema-value alignment \\
        OFT (post-finetune) & 27.35 & 100.00 & Large gain in normalized exact match; syntax remains robust \\
        Delta (OFT - Baseline) & +27.20 & 0.00 & Major semantic improvement at constant parse success \\
        \bottomrule
    \end{tabularx}
\end{table}

\insertfigure{../outputs/plots/eval_metrics.png}{0.75}{Bar chart of evaluation metrics (optional if table is sufficient).}

\subsection{Qualitative Examples}
\begin{table}[H]
    \centering
    \caption{Representative Before/After Predictions}
    \begin{tabularx}{\linewidth}{>{\raggedright\arraybackslash}p{0.19\linewidth}X}
        \toprule
        Field & Content \\
        \midrule
        Question & What's the name of the carabins stadium? \\
        Schema snippet & \texttt{CREATE TABLE table\_name\_55 (football\_stadium VARCHAR, team VARCHAR)} \\
        Gold SQL & \texttt{SELECT football\_stadium FROM table\_name\_55 WHERE team = "carabins"} \\
        Baseline prediction & \texttt{SELECT football\_stadium FROM table\_name\_55 WHERE team = 'carabins' LIMIT 1;} \\
        OFT prediction & \texttt{SELECT football\_stadium FROM table\_name\_55 WHERE team = "carabins"} \\
        Notes & OFT output exactly matches normalized gold SQL, while baseline adds unnecessary \texttt{LIMIT 1}. \\
        \bottomrule
    \end{tabularx}
\end{table}

\noindent Additional qualitative examples are available in \texttt{outputs/eval/qualitative\_examples.json}.



\end{document}
